\section{Discussion}
\label{sec:discussion}

\paragraph{Two complementary failure modes.}
The two studies reveal two complementary ways in which multi-turn context can
degrade reasoning.

\begin{itemize}
  \item \textbf{Short tasks (Study 1).} Retaining the assistant's own replies
    lowered final success relative to an information-matched context without
    them. This replicates the direction reported by \citet{huang2026benefit}
    under a frozen, task-paired analysis. Mistral-7B illustrates one mechanism:
    with its own replies in context it tended to remain in an explicit ``need
    more information'' state, while the same model given the same information
    without those replies attempted an answer.
  \item \textbf{Long horizons (Study 2).} With a steadily updated state, accuracy
    collapsed within about twenty turns, even at fixed difficulty. A verified
    state summary avoided the collapse \pending{in all three models}.
\end{itemize}

The studies differ in tasks, models and controls. We therefore do not claim that
the short-task effect causally grows into the long-horizon collapse. Study 2 also
cannot separate three possible causes: history length, the number of state
updates, and error accumulation.

\paragraph{Why quality metrics stay silent.}
Consistency, coherence and stability measure local properties of responses. A
model that has committed to a wrong premise or a wrong state can reason
coherently and repeatably from it. In our replay these metrics even rose as
accuracy fell, which suggests that accumulated context makes responses more
anchored.

The paired process measure fared no better than the current readiness state
itself. Accuracy-based detectors, for their part, are limited by the noise of a
single stream. Monitoring designs that assume quality metrics fall before
accuracy are therefore unlikely to work in this setting. Token-level uncertainty
signals such as those used by ERGO \citep{khalid2025ergo} are a different class
of signal. We did not test them, because our protocols did not collect token
probabilities.

\paragraph{Implications for practice.}
Three practices follow from our results:
\begin{itemize}
  \item Do not let the raw history grow indefinitely in tasks that maintain
    state. Re-ground the model periodically in a compact, verified state
    representation.
  \item Monitor by comparison: against a verified-state or history-free
    re-query, or across many concurrent sessions. Single-stream quality
    trajectories are not a substitute.
  \item Before using a generic drift detector as a comparator, verify that it
    can fire on streams of the intended length.
\end{itemize}

\paragraph{Limitations.}
\begin{itemize}
  \item \textbf{Scope.} Both studies use arithmetic reasoning with exact answers
    and models available through one academic gateway, and the two studies use
    different models. Effects may differ for other domains, newer models or other
    serving stacks.
  \item \textbf{Net effects.} Study 1 estimates the net effect of retaining
    assistant replies, which contain both harmful commitments and potentially
    useful work. Study 2's control is a packaged intervention, not a mechanism
    test.
  \item \textbf{Observability.} Accuracy, robustness and efficiency were
    unobservable before the final turn in Study 1. Their early-warning value
    there was not tested.
  \item \textbf{Exploratory analyses.} The Study 2 monitoring analyses, the CRL
    persistence check, the accuracy-loss target and the stream replay are
    exploratory or post-hoc, and are labelled as such.
  \item \textbf{Study 2 acquisition.} A funding pause staged the acquisition.
    \pending{The frozen completion of the control arm, and the status of the
    remaining consistency, robustness and stability samples, which are reported
    for GPT-4o-mini only.}
  \item \textbf{Single investigator.} The audit of Study 1 tasks and the
    calibration of Study 2 were carried out by one investigator.
\end{itemize}
